
Pour essayer de comprendre le comportement du pendule en regime libre du balencier spiral, nous utilisons l'application Phidgets MATLAB pour mesure l'angle du balancier.
\\
Nous complementons les mesures avec l'utilisation d'un amperemetre pour mesurer le courant induit dans le frein. Ainsi nous tentons d'extraire les rapports entre l'angle et les courants pour 
\\
%Definir la largeur : \includegraphics[width=0.7\textwidth]{nom_image.png}.
%Définir la hauteur : \includegraphics[height=8cm]{nom_image.png}.
%Redimensionner proportionnellement : \includegraphics[scale=0.75]{nom_image.png}.

\begin{figure}[ht]
\centering
\includegraphics[scale=0.5]{schema_force.png}
\caption{Schema electrique pour la prise de mesure en regime d'oscillations libres}
\label{fig:samplesetup}
\end{figure}

\\
Avant de mesurer les oscillations avec un courant dans le frein. 
\\
\begin{figure}[ht]
\centering
\includegraphics[scale=0.5]{schema_force.png}
\caption{Schema electrique pour la prise de mesure en regime d'oscillations forcees}
\label{fig:samplesetup}
\end{figure}

